\scanpage{76}
\mspage{12}
Au niveau des multiplicités modulaires \inserted{(analytiques réelles)},
on trouve sûrement un morphisme pour
celles-ci
\[
\begin{aligned}
&\deleted{MB_{g,\nu}\unsure{-|\to}}\\
&MDB_{g,\nu}\longrightarrow M_{\mathrm{st}}N\deleted{{}_{g,\nu}}D_{g,\nu}
\end{aligned}
\]
où $MDB_{g,\nu}$ désigne la multiplicité modulaire
($M$) pour les surfaces holomorphes compactes
``découpées'' ($D$) à bord ($B$) de type $g,\nu$, et
$M_{\mathrm{st}}ND_{g,\nu}$ celle pour les courbes algébriques
stables (de type loc. constant) \inserted{de type $g,\nu$} à structure de
propagation normale ($N$) et à système \inserted{adm} de
disques ($D$) autour des pts marqués.

Cette flèche est compatible avec les
flèches canoniques vers la cat. des
types combinatoires \inserted{-toriques} de tels objets
(diagrammes $MD$ (Mumford-Deligne ou
``modulaires'') (\deleted{avec pond} \inserted{incluant} une pondération
des sommets) à structures toriques
associées aux arêtes\ldots{} donnant lieu un
diagramme commutatif
% archive sheet number: 78
